% !TEX root = ../ts_notes.tex
\documentclass[12pt,fleqn]{article}
\usepackage{../vkCourseML}
\newcounter{chapter}
\hypersetup{unicode=true}
\interfootnotelinepenalty=10000
\newcommand{\dx}[1]{\,\mathrm{d}#1} % маленький отступ и прямая d
\begin{document}
\section{Компоненты временного ряда}


Самый простой способ разложить временной ряд - это использовать предпосылку об аддитивности.
Тогда ряд можно представить в виде суммы нескольких компонент:

\[
y_t = t_t + s_t + c_t + e_t,
\]

где \(t_t\) - тренд, \(s_t\) - сезонность, \(c_t\) - цикличность, а \(e_t\) - ошибка или остаток.
На практике не всегда выделяют все четыре компоненты; часто используют более простую схему

\[
y_t = t_t + s_t + e_t.
\]

Тренд - это медленно меняющийся уровень ряда.
Это определение не является строгим: разные алгоритмы могут по-разному отделять тренд от остальных компонент.

Сезонность - это повторяющееся колебание с фиксированным периодом.
Например, недельная сезонность в дневных продажах означает, что похожий паттерн повторяется каждые семь дней.

Цикличность - это колебание без фиксированного периода.
Типичный пример - экономические циклы: рост, охлаждение, спад и восстановление могут повторяться, но длина цикла заранее неизвестна.

Остаток - это то, что осталось после удаления объяснимых компонент.
Если в остатках сохраняется автокорреляция, это сигнал, что модель не захватила какую-то структуру ряда.

Декомпозиция полезна для прогноза.
Тренд можно продолжить простой моделью, сезонность можно повторять или моделировать отдельно, а остаток либо не прогнозировать, либо подавать в более гибкую модель вместе с внешними факторами.
Например, для числа игроков в игре можно отдельно описать тренд после релиза, суточную сезонность, а остаток объяснять событиями вроде сбоев сервера или внешних шоков.


\section{Алгоритмы сглаживания}


STL и MSTL - алгоритмы разложения временного ряда на компоненты.
STL выделяет сезонность, тренд и остаток.
MSTL обобщает эту идею на несколько сезонностей, что важно, например, для дневных данных с недельной и годовой сезонностью.

Чтобы понять эти алгоритмы, нужно сначала обсудить сглаживание: как из шумного ряда получить более гладкую последовательность

\[
(y_t)\longrightarrow(\tilde y_t).
\]


\subsection{Moving average}


Самый простой алгоритм сглаживания - скользящее среднее.
Для окна ширины \(5\) можно записать

\[
\tilde y_t
=
\frac{
y_{t-2}+y_{t-1}+y_t+y_{t+1}+y_{t+2}
}{5}.
\]

Чем шире окно, тем сильнее сглаживание.
Узкое окно сохраняет больше локальных колебаний, широкое сильнее убирает шум, но может стереть полезную структуру.

У скользящего среднего есть проблема на краях ряда.
Для первых и последних наблюдений не хватает соседей слева или справа.
Возможные решения:

\begin{itemize}
  \item сократить длину итогового сглаженного ряда;
  \item экстраполировать ряд за границы;
  \item заполнить недостающие значения последним доступным значением;
  \item использовать асимметричные окна.
\end{itemize}

Универсального ответа нет: выбор зависит от частоты данных, длины ряда и цели сглаживания.


\subsection{OLS}


Еще один способ сгладить ряд - построить регрессию на временную шкалу.
Пусть \(x_t\) обозначает время.
Тогда можно оценить модель

\[
y_t = \beta_1+\beta_2 x_t+\varepsilon_t
\]

и взять сглаженное значение

\[
\tilde y_t=\hat\beta_1+\hat\beta_2 x_t.
\]

Такая модель часто слишком груба: она описывает только глобальный линейный тренд.
Но она важна как переход к локальным регрессиям.


\subsection{LOESS}


LOESS означает local regression, а LOWESS - locally weighted regression.
Идея состоит в том, чтобы для каждой точки строить локальную взвешенную регрессию: ближайшие наблюдения получают больший вес, дальние - меньший.

Для текущей точки \(x\) можно записать функционал

\[
Q(\hat\beta_1,\hat\beta_2;x)
=
\sum_{t=1}^T
K(x_t,x)
\left(
y_t-(\hat\beta_1+\hat\beta_2 x_t)
\right)^2.
\]

Здесь \(K(x_t,x)\) - ядро, или функция весов.
Она определяет, насколько сильно наблюдение \(x_t\) влияет на локальную регрессию около точки \(x\).

Если взять индикаторное ядро

\[
K(x_t,x)=
\begin{cases}
1, & x_t \text{ входит в } h \text{ ближайших наблюдений к } x,\\
0, & \text{иначе},
\end{cases}
\]

то получится метод, близкий по духу к скользящему окну.
Чаще используют более плавные ядра, например экспоненциальное:

\[
K(x_t,x)
=
\exp\left(
-\frac{(x_t-x)^2}{h^2}
\right),
\qquad h>0.
\]

Параметр \(h\) управляет чувствительностью.
При большом \(h\) веса убывают медленно, и метод приближается к глобальной регрессии.
При малом \(h\) используются почти только ближайшие точки, и сглаживание становится очень локальным.

После оценки локальной регрессии сглаженное значение в точке \(x_t\) записывается как

\[
\tilde y_t=\hat\beta_1+\hat\beta_2 x_t.
\]

LOESS/LOWESS нужен как технический аппарат для дальнейшего обсуждения STL и MSTL: эти алгоритмы используют локальное сглаживание, чтобы отделять тренд и сезонные компоненты от остатка.

\section{STL}


Лекция посвящена декомпозиции временного ряда на компоненты и двум алгоритмам из семейства STL.
Сначала разбирается базовый STL, который работает с одной сезонной частотой, затем его обобщение MSTL для рядов с несколькими сезонностями.

Основная задача декомпозиции - представить наблюдаемый ряд как сумму более простых компонент:

\[
y_t=t_t+s_t+e_t.
\]

Здесь \(t_t\) - плавно меняющийся тренд, \(s_t\) - сезонная компонента с жесткой периодичностью, а \(e_t\) - остаток, то есть все, что не удалось объяснить трендом и сезонностью.
Циклическую компоненту в базовом STL отдельно не выделяют: цикличность встречается реже, ее сложнее устойчиво оценивать, и для нее обычно нужны более специальные модели.

\begin{figure}[htbp]
  \centering
  \fbox{\parbox{0.82\linewidth}{\centering Заготовка для рисунка: Схематическое аддитивное разложение STL}}
  \caption{Рисунок 1.
  Схематическая иллюстрация аддитивного разложения ряда на тренд, сезонность и остаток.}
  \label{fig:lecture_02_stl_components}
\end{figure}


\subsection{Идея STL}


STL - это эвристический, но очень практичный алгоритм декомпозиции.
Он не дает единственно "правильного" разложения ряда: разные настройки сглаживания тренда и сезонности могут давать разные, но все еще осмысленные декомпозиции.
Поэтому STL стоит воспринимать как хорошо подобранную последовательность сглаживаний, а не как строго выведенную теорему.

Сильные стороны STL:

\begin{itemize}
  \item он давно используется и хорошо проверен на практике;
  \item он достаточно робастен к разным формам данных;
  \item его основные шаги построены из уже знакомых операций: разбиения ряда по сезонным позициям, LOESS-сглаживания и скользящих средних;
  \item во многих библиотеках уже подобраны адекватные параметры по умолчанию.
\end{itemize}

Главное ограничение базового STL: он работает только с одной сезонной частотой.
Например, для месячных данных естественный период сезонности часто равен \(12\), а для часовых данных может быть период \(24\), если нужно выделить суточный паттерн.


\subsection{Компоненты и начальная инициализация}


В алгоритме будем использовать три компоненты:

\[
s_t,\qquad
t_t,\qquad
e_t.
\]

На первом шаге инициализация задается просто:

\[
s_t := y_t,\qquad
t_t := 0,\qquad
e_t := 0.
\]

Такой старт выглядит грубо, но дальше компоненты итеративно уточняются.
Важная идея состоит в том, что тренд и сезонность оцениваются не один раз, а постепенно очищаются друг от друга.


\subsection{Алгоритм STL}


Пусть выбран сезонный период \(m\).
Для месячных данных с годовой сезонностью обычно \(m=12\).
Тогда шаги STL можно записать так.

\begin{figure}[htbp]
  \centering
  \fbox{\parbox{0.82\linewidth}{\centering Заготовка для рисунка: Схема шагов STL}}
  \caption{Рисунок 2.
  Восстановленная по рукописному PDF схема шагов STL: разбиение по периоду, локальные сглаживания, выделение сезонности, тренда и остатка.}
  \label{fig:lecture_02_stl_flow}
\end{figure}

\begin{enumerate}
  \item Инициализировать компоненты:

\[
s_t := y_t,\qquad
t_t := 0,\qquad
e_t := 0.
\]

  \item Удалить текущую оценку тренда:

\[
y_t^{\mathrm{det}}
=
y_t-t_t.
\]

На первой итерации этот шаг почти ничего не делает, потому что \(t_t=0\).
Но при повторных проходах в тренде уже будет лежать предыдущая оценка.

  \item Разрезать детрендированный ряд по сезонным позициям.
  Если \(m=12\), то ряд раскладывается на 12 подрядов: все январи, все феврали, ..., все декабри.
  В общем виде это можно понимать как операцию

\[
y_t^{\mathrm{det}}
\longrightarrow
\left(y_t^{(1)},y_t^{(2)},\ldots,y_t^{(m)}\right).
\]

Смысл простой: чтобы оценить, например, январский сезонный эффект, нужно смотреть только на наблюдения, попавшие в январскую позицию разных лет.

  \item Сгладить каждый сезонный под ряд с помощью LOESS.
  Для каждой сезонной позиции получается сглаженная версия:

\[
C_t^{(j)}=\mathrm{LOESS}\left(y_t^{(j)}\right),
\qquad j=1,\ldots,m.
\]

Здесь появляются первые важные настройки: окно сглаживания, ядро локальной регрессии и другие параметры LOESS.
Чем сильнее сглаживание, тем более стабильной становится сезонная компонента, но тем больше риск потерять реальные локальные изменения.

  \item Собрать сглаженные подряды обратно в один ряд \(C_t\), поставив элементы в исходный временной порядок:

\[
\left(C_t^{(1)},\ldots,C_t^{(m)}\right)
\longrightarrow
C_t.
\]

  \item Сильно сгладить \(C_t\), чтобы получить прокси для тренда внутри сезонной оценки.
  На слайде этот шаг описан как комбинация скользящих средних и LOESS:

\[
L_t
=
\mathrm{LOESS}\left(
\mathrm{MA}\left(
\mathrm{MA}\left(C_t\right)
\right)
\right).
\]

Точная техническая реализация может отличаться между библиотеками.
Интуитивно \(L_t\) - это очень гладкая часть \(C_t\), которую не хочется считать сезонностью.

  \item Выделить сезонность:

\[
s_t=C_t-L_t.
\]

То есть из предварительной сглаженной сезонной кривой вычитается слишком гладкая часть, похожая на тренд.
В результате остается сезонный паттерн.

  \item Обновить тренд, удалив из исходного ряда найденную сезонность:

\[
t_t
=
\mathrm{LOESS}\left(y_t-s_t\right).
\]

Этот шаг нужен, чтобы тренд оценивался уже по ряду, очищенному от сезонности.

  \item Повторить шаги 2-8 несколько раз.
  Обычно большого числа итераций не нужно: после нескольких проходов изменения в тренде и сезонности становятся маленькими.

  \item После финального обновления вычислить остаток:

\[
e_t
=
y_t-t_t-s_t.
\]

Итоговое разложение остается аддитивным:

\[
y_t
=
t_t+s_t+e_t.
\]

Важно помнить, что аддитивная декомпозиция подходит не всегда.
Если амплитуда сезонности зависит от уровня тренда, может понадобиться мультипликативная модель.
Такие случаи позже появляются в ETS-моделях.

\end{enumerate}

\subsection{Как выбирать период}


Для STL нужно заранее указать период сезонности.
В разных библиотеках слова "частота" и "период" могут использоваться немного по-разному, поэтому полезно держать в голове практический смысл: период - это число наблюдений между повторениями сезонного паттерна.

Примеры:

\begin{itemize}
  \item месячные данные с годовой сезонностью: \(m=12\);
  \item дневные данные с недельной сезонностью: \(m=7\);
  \item часовые данные с суточной сезонностью: \(m=24\);
  \item часовые данные с недельной сезонностью: \(m=24\cdot 7\).
\end{itemize}

Период можно искать несколькими способами:

\begin{itemize}
  \item визуально по графику ряда;
  \item по графику автокорреляции;
  \item по предметному знанию о данных;
  \item через частотные методы вроде преобразования Фурье, если они уместны.
\end{itemize}

Автокорреляция на лаге \(k\) - это корреляция между \(y_t\) и \(y_{t-k}\).
Если, например, на лаге \(12\) виден устойчивый пик, это сигнал в пользу сезонности с периодом \(12\).
Частная автокорреляция похожа по духу, но измеряет связь между \(y_t\) и \(y_{t-k}\) после учета промежуточных лагов \(y_{t-1},\ldots,y_{t-k+1}\).

При выборе периодов действует практическое правило: если более простой период уже дает хорошую декомпозицию, не стоит без необходимости добавлять много сезонных компонент.
Каждая дополнительная компонента усложняет последующее прогнозирование: нужно будет прогнозировать больше частей, и ошибка может накапливаться.


\subsection{Остатки и качество разложения}


Хорошее разложение должно оставлять в остатках как можно меньше структуры.
Если после удаления тренда и сезонности остатки все еще имеют выраженную автокорреляцию, значит в ряде осталась неучтенная закономерность.

В идеале остатки должны быть похожи на белый шум: у них нет линейной автокорреляции, а оставшиеся колебания выглядят как непредсказуемая ошибка.
На практике это проверяют визуально, через ACF-графики и позже через статистические тесты.

Если остатки не похожи на белый шум, это не обязательно означает, что STL бесполезен.
Разложение все еще может быть полезным для прогноза:

\begin{itemize}
  \item тренд можно прогнозировать простой моделью или сглаживанием;
  \item сезонность можно повторять по прошлому периоду или усреднять по прошлым периодам;
  \item остатки можно моделировать отдельно более сложной моделью.
\end{itemize}

Именно поэтому декомпозиция часто используется как промежуточный шаг в построении прогноза.


\subsection{Робастность и структурные разрывы}


В реализациях STL обычно есть робастный вариант, который аккуратнее относится к выбросам.
Это важно, если в ряде есть резкие провалы или всплески: авария, шок спроса, обновление продукта, ошибка измерения.
Обычный STL может частично размазать такой выброс по тренду или сезонности.

Со структурными разрывами нужно быть особенно осторожным.
Если ряд после некоторого момента стал генерироваться другим процессом, есть несколько вариантов:

\begin{itemize}
  \item попытаться устранить разрыв вручную, если причина известна и есть понятная корректировка;
  \item строить отдельную модель только на актуальном фрагменте после разрыва;
  \item добавить индикатор режима или другую объясняющую переменную;
  \item использовать модель, которая явно учитывает смену режима.
\end{itemize}

Для STL в такой ситуации часто разумно либо предварительно поправить разрыв, либо оценивать декомпозицию только на последнем устойчивом участке.
Но перед этим нужно понять, действительно ли разрыв означает смену параметров процесса.


\subsection{Параметры STL в StatsModels}


В \texttt{statsmodels} для STL особенно важны первые параметры:

\begin{itemize}
  \item \texttt{period} - сезонный период, например \(12\) для месячных данных;
  \item параметры сглаживания сезонности;
  \item параметры сглаживания тренда;
  \item \texttt{robust} - включение робастной версии для выбросов.
\end{itemize}

Чем больше окно сглаживания сезонности, тем более гладкой получается сезонная компонента.
Чем больше окно сглаживания тренда, тем более гладким становится тренд.
Если сгладить тренд слишком сильно, остатки могут начать содержать очевидную структуру: тренд окажется слишком грубым и не заберет часть реальной динамики.

В MSTL похожая настройка называется \texttt{windows}: для каждого периода можно задать свое окно сезонного сглаживания.
Обычно эти окна должны быть нечетными.


\subsection{Пример с месячными данными}


Для месячного ряда производства или потребления электроэнергии естественный первый кандидат на сезонность - период \(12\).
На графике ряда годовой паттерн обычно виден напрямую, а на ACF-графике появляются пики на лагах \(12,24,\ldots\).

Иногда можно увидеть и пики на лагах \(3\) или \(6\).
Но они могут быть вложены в годовую сезонность: если годовой паттерн сложный, внутри него могут возникать квартальные или полугодовые совпадения.
Поэтому не стоит автоматически добавлять периоды \(3\), \(6\) и \(12\) одновременно.
Сначала нужно проверить, хватает ли периода \(12\).
Если разложение выглядит хорошо, более простая модель предпочтительнее.

При этом остатки после STL могут все равно содержать зависимость.
Тогда декомпозиция полезна, но не завершает моделирование: тренд, сезонность и остаток придется прогнозировать или моделировать отдельно.


\section{MSTL}


MSTL расшифровывается как multi-seasonal trend decomposition using LOESS.
Это обобщение STL на случай нескольких сезонностей.
Например, в дневных данных могут одновременно присутствовать недельная и месячная сезонность, а в часовых данных - суточная и недельная.

Вместо одной сезонной компоненты получаем несколько:

\[
y_t
=
t_t
+s_t^{(1)}
+s_t^{(2)}
+e_t.
\]

Для большего числа сезонностей формула расширяется очевидным образом:

\[
y_t
=
t_t
+\sum_{j=1}^{K}s_t^{(j)}
+e_t.
\]

Здесь каждая \(s_t^{(j)}\) соответствует своему периоду.

\begin{figure}[htbp]
  \centering
  \fbox{\parbox{0.82\linewidth}{\centering Заготовка для рисунка: Схема MSTL}}
  \caption{Рисунок 3.
  Восстановленная схема MSTL: сезонности выделяются последовательно, затем уточняются коррекционными проходами.}
  \label{fig:lecture_02_mstl_flow}
\end{figure}


\subsection{Почему начинают с высоких частот}


В MSTL сезонности обычно достают последовательно от более высокой частоты к более низкой.
Если говорить через периоды, то сначала берут меньший период.
Например, в дневных данных с недельной и месячной сезонностью сначала берут \(7\), потом около \(30\).

Это эмпирическое правило, не теорема.
Интуиция такая: за фиксированный промежуток времени более высокая частота повторяется больше раз, поэтому ее можно оценить точнее.
Если есть один год дневных данных, недельных циклов в нем много, а месячных намного меньше.

Еще одна причина - наложение частот.
Например, конец недели может совпасть с концом месяца.
Тогда первый проход STL может забрать часть эффекта не полностью или смешать его с другой сезонностью.
Поэтому MSTL делает последовательные проходы и затем корректирует уже найденные компоненты.


\subsection{Алгоритм MSTL для двух сезонностей}


Пусть есть две сезонности: первая с более высокой частотой, вторая с более низкой.
Обозначим их \(s_t^{(1)}\) и \(s_t^{(2)}\).

\begin{enumerate}
  \item Применяем STL к исходному ряду с первым периодом и достаем первую сезонность:

\[
y_t
=
s_t^{(1)}
+y_t^{(-1)}.
\]

Здесь \(y_t^{(-1)}\) - глобальный остаток после удаления первой сезонности.

  \item Применяем STL ко второму ряду \(y_t^{(-1)}\), уже с низкой частотой:

\[
y_t^{(-1)}
=
s_t^{(2)}
+y_t^{(-1,2)}.
\]

Теперь \(y_t^{(-1,2)}\) - остаток после удаления двух сезонностей.
Если сезонностей больше, этот шаг повторяется для остальных периодов.

  \item Делаем корректировку.
  Возвращаем первую сезонность к текущему остатку и снова раскладываем STL с первым периодом:

\[
y_t^{(-1,2)}
+s_t^{(1)}
=
s_t^{(1)*}
+y_t^{*(-1,2)}.
\]

Так первая сезонность уточняется с учетом того, что часть низкочастотной структуры уже была удалена.

  \item Аналогично уточняем вторую сезонность.
  После последовательной корректировки получаем финальные компоненты:

\[
y_t
=
s_t^{(1)*}
+s_t^{(2)*}
+t_t^{**}
+e_t^{**}.
\]

Эту процедуру можно повторять несколько раз, но на практике после первых проходов корректировки обычно становятся маленькими.
Поэтому бесконечно крутить цикл нет смысла.

\end{enumerate}

\subsection{Примеры MSTL}


В синтетическом примере можно сгенерировать часовой ряд с ломаным трендом, суточной сезонностью \(24\) и недельной сезонностью \(24\cdot 7\).
Такая структура похожа на онлайн-игру или сервис: внутри дня есть пики активности, а внутри недели есть отдельный рабоче-выходной паттерн.
При правильно указанных периодах MSTL хорошо выделяет обе сезонности и тренд.

\begin{figure}[htbp]
  \centering
  \fbox{\parbox{0.82\linewidth}{\centering Заготовка для рисунка: Схематический пример ряда с двумя сезонностями}}
  \caption{Рисунок 4.
  Схематическая иллюстрация часового ряда с суточной и недельной сезонностью.
  Это не исходные данные лекции, а восстановленная по обсуждению качественная схема.}
  \label{fig:lecture_02_multiple_seasonality}
\end{figure}

В более реальном примере с часовым потреблением электроэнергии в Австралии видны:

\begin{itemize}
  \item суточная сезонность: днем потребление выше, ночью ниже;
  \item недельная сезонность: в рабочие дни потребление выше, на выходных ниже;
  \item остаточная структура, которую две сезонности не объясняют полностью.
\end{itemize}

Если в остатках остаются повторяющиеся пики, нужно разбираться дальше.
Возможные причины:

\begin{itemize}
  \item пропущена еще одна сезонность;
  \item выбранные окна сглаживания слишком сильные или слишком слабые;
  \item в данных есть структурные события;
  \item остаток зависит от внешних факторов, которые не попали в модель.
\end{itemize}

В таких случаях помогает ACF-график остатков и предметный анализ конкретных дат.
Иногда нужно добавить новую сезонность, иногда - инженерить признаки и строить отдельную модель для остатка.


\subsection{Практический итог}


STL и MSTL полезны как базовый инструмент декомпозиции:

\begin{itemize}
  \item STL подходит, когда есть одна главная сезонность;
  \item MSTL подходит, когда сезонностей несколько;
  \item периоды нужно задавать заранее и проверять визуально или через автокорреляции;
  \item остатки после декомпозиции нужно анализировать отдельно;
  \item робастный режим полезен при выбросах;
  \item структурные разрывы лучше обрабатывать до декомпозиции или моделировать явно.
\end{itemize}

Главная мысль лекции: декомпозиция не является полностью автоматическим доказательным разложением ряда на "истинные" компоненты.
Это практическая процедура сглаживания и фильтрации, которая помогает получить интерпретируемые части ряда и упростить последующее прогнозирование.
\newpage
% \begin{thebibliography}{0}
% \end{thebibliography}

\end{document}
